% STATUS: version with results (2026-10-10). Numbers: Project/docs/checkpoint_2026-10-10.md.

La retrieval-augmented generation (RAG) è ormai il modo standard per dare ai modelli linguistici accesso a documenti
privati, e oggi l'intera pipeline può essere eseguita sul dispositivo dell'utente con hardware di fascia consumer, a
vantaggio della privacy, dell'uso offline e dei costi. Il recupero dei documenti, tuttavia, non elimina le allucinazioni:
le risposte generate contengono ancora frasi che il contesto recuperato non supporta o che lo contraddicono. Individuarle
è un prerequisito per assistenti locali affidabili e, poiché il controllo va eseguito su ogni risposta, deve trovare
posto sulla stessa macchina, dalle risorse limitate, che ospita il generatore.

La verifica segue oggi due paradigmi. I verificatori basati su encoder, come AlignScore, MiniCheck e HHEM, sono piccoli e
veloci, ma restituiscono un unico punteggio fisso per ogni coppia di documento e frase. I giudici generativi sono
flessibili, ma scrivono i verdetti token per token, per cui il loro costo cresce con la quantità di testo prodotto; i
giudici piccoli violano spesso il formato di output, e la confidenza che esprimono è raramente utilizzabile come
probabilità. Di recente è emersa una terza famiglia: i \textbf{modelli System One}, come il proprietario Jev e i modelli
aperti SemIf e Rizzo-Flow, che rispondono a domande chiuse in un unico forward pass leggendo dai logit la probabilità di
ciascuna opzione di risposta. Le evidenze a loro favore provengono per lo più dagli stessi autori, e le implementazioni
aperte sono generaliste. Non era ancora stato studiato in modo controllato se un System One compatto, specializzato in un
unico task, possa eguagliare i giudici generativi e gli encoder dedicati, da dove derivi il suo vantaggio e se le sue
probabilità siano utilizzabili, il tutto su un portatile.

Questa tesi progetta, addestra e valuta \textbf{GroundOne}, un System One specializzato in un'unica decisione: stabilire,
per ogni frase di una risposta, se il documento recuperato la supporta, la contraddice o non la menziona. Ogni frase
della risposta è seguita da un marker, e i verdetti di tutte le $K$ frasi si leggono dai logit di tre lettere di risposta
in un unico forward pass, senza generare testo. Modelli Qwen2.5 da 0.5B e 3B parametri sono sottoposti a fine-tuning con
LoRA sulle risposte di RAGTruth annotate manualmente, con una loss pesata per classe e ristretta alle tre lettere; una
temperatura e un bias stimati sul dev set calibrano le probabilità, e una soglia scelta sugli stessi dati le traduce in
etichette coerenti e in segnalazioni secondo la policy strict o lenient. Tutto, dall'addestramento alla messa in
servizio, gira con MLX su un portatile Apple M4 con 16\,GB di memoria. La valutazione usa tre test set costruiti e
congelati prima dell'addestramento (la partizione di test di RAGTruth, LLM-AggreFact e VitaminC) e confronta GroundOne con
indizi che non leggono mai il documento, con i modelli non addestrati, con giudici generativi della stessa dimensione, con
e senza probabilità, con un System One generalista e con encoder specializzati, riportando intervalli calcolati con un
bootstrap a cluster.

Sul test set di RAGTruth GroundOne raggiunge una balanced accuracy a tre classi di 0.769 con 0.5B parametri e di 0.832
con 3B, contro 0.37 e 0.44 dei modelli non addestrati; la sua AUROC (0.905 e 0.940) supera quella degli encoder (circa
0.82), e il modello supera nettamente sia un System One generalista otto volte più grande sia un decision model
generalista della stessa dimensione (0.657), che a sua volta è molto più forte sul test fuori dominio VitaminC. Su
LLM-AggreFact il 3B è alla pari con MiniCheck-RoBERTa (74.8 contro 74.0, differenza non significativa), sopra il modello
7B zero-shot pubblicato e sotto il System One chiuso Jev, misurato in modo indipendente (78.6). Per risposte di
quattordici o più frasi GroundOne è circa due volte più veloce di un giudice sintetico e sei volte più veloce di uno
verboso della stessa dimensione, con un output sempre valido; consuma inoltre meno energia (5.5 e 33 joule per risposta),
sebbene con un fattore inferiore a quello osservato sul tempo, perché un prefill lungo assorbe più potenza del decoding.
Quando si chiede loro lo stesso output, i giudici generativi forniscono probabilità che distinguono a stento le decisioni
corrette da quelle errate, mentre la calibrazione porta l'Expected Calibration Error di GroundOne sul test set a 0.019. A
parità di pesi, restringere la decisione alle tre lettere è più accurato che prendere il token che il modello
genererebbe, di 0.25 per il 0.5B e di 0.07 per il 3B. I limiti sono altrettanto chiari: il trasferimento a un dominio
diverso è debole, le allucinazioni nascoste in frasi che per il resto ricalcano la fonte restano difficili da
individuare, e la ricetta di addestramento non è del tutto robusta rispetto alla scelta del seed.
